\chapter{Research methodology}
\label{chapter:methodology}

The approach is design-first and simulation-first. Candidate concepts are traded off, then the selected design is advanced and analysed with thermomechanical and structural simulation before any physical prototyping. PLASMO is currently a feasibility study, so experimental work stays at conceptual or proof-of-concept level.

The methodology carries the consequence of \Cref{sec:review-method}. Because the published literature cannot be compared like for like, the comparison that decides between the carried architectures has to be generated within the thesis. All candidates are modelled parametrically on one common set of assumptions, in one material class and sized on one common basis, so that the differences between them follow from their architecture rather than from the conditions under which each was built. This is the one like-for-like comparison the thesis can make, and it is why the trade-off rests on the modelling of \Cref{sec:phase1} rather than on the literature.

The methodology also follows the sequence that recurs across the predecessor theses in the Deployable Space Telescope line: requirements stated early and numbered; an option survey with explicit early elimination; criteria defined before concepts are scored, followed by a formal trade-off; detailed design of the selected concept; ESATAN thermal analysis together with finite-element modal and thermoelastic analysis, hand-checked, with a formal error budget where accuracy matters; and conclusions that answer the research questions in order. ECSS handbooks are used as a source of trade-off methodology rather than as strict requirements.

\section{Phase 1, to the midterm: from the carried architectures to one design}
\label{sec:phase1}
\phasestart{1}

\activity{criteria}{define the trade-off criteria} The criteria and their weights are fixed before any candidate is modelled or scored. They are derived from the architectural properties of \Cref{sec:properties}, and they have to span the dimensions on which \Cref{rq:concepts} asks the architectures to be compared: volumetric packing ratio, the enclosure volume each one forces, bare deployment accuracy, in-orbit stability, stiffness, mass and complexity. Required active correction is carried as a candidate criterion rather than a settled one: it discriminates only if an architecture rules a compensation scheme in or out, and it falls away entirely if every carried architecture leaves every scheme available. Defining the set first is what determines which quantities the models have to produce.

% \activity{inputs}{define the model inputs and outputs} Before anything is built, the parameter set is written down in full: which quantities are inputs, which are swept, and which are held fixed so that the architectures can be compared fairly. Two of these need a decision rather than a value.

% \begin{itemize}
%     \item \textbf{Stiffness} is the binding one. It has to be held comparable across architectures, because it drives the dimensioning of the members and through them the mass and the stowed volume, yet a wireframe of lines and points has no cross-section, so nothing in the geometry fixes a wall thickness. A single assumed thickness for every family is not a fair basis: trusses and masts are not monolithic members and are typically larger but lighter; the driven linkage has members of different lengths and is only comparable once fully deployed; telescoping booms have progressively smaller segments, so their stiffness is not simple to approach. For that last case the telescopic-boom work of Wilting gives worked formulas that can be adopted rather than re-derived.
%     \item \textbf{The packing ratio boundary.} The figure is the volumetric ratio of the M2 mechanism stowed to the same mechanism deployed, bounded where it mounts to the rest of the spacecraft and running to its furthest extent: not the individual boom, not the satellite, and not other deployables. Each architecture reports two envelopes, the minimum envelope around the wireframe and the cylinder or box it fits into, giving a range rather than one number, with a qualitative note on how practical the resulting shape is. The volume the baffle is forced to occupy is recorded as a separate figure alongside the ratio, never inside it, because a concept that flares at its base when stowed or that swings outside the bus pushes the baffle outward and so enlarges the stowed volume at system level.
% \end{itemize}

\activity{inputs}{define the model inputs and outputs} Before anything is built, the parameter set is written down in full: which quantities are inputs, which are swept, which are held fixed so that the architectures can be compared fairly, and which are deliberately left out. Four of these need a decision rather than a value.

\begin{itemize}
    \item \textbf{Member sizing and stiffness.} Stiffness is the binding parameter, because it drives the dimensioning of the members and through them the mass and the stowed volume, yet a wireframe of lines and points has no cross-section, so nothing in the geometry fixes a wall thickness. The sizing is therefore derived rather than assumed. Each architecture is built from wireframe members in one common material, loaded by a notional reference force at M2, and the member thickness is back-calculated from that load, so that every architecture is sized to the same structural performance over the same deployment length. Mass and stowed volume then follow from the sizing instead of being set by hand, and the trade between many slender members and few stiff ones becomes visible in the result. A starting thickness is taken from the members of the earlier PLASMO boom design and scaled with length. The quantity that drives the sizing is the tip deflection under the reference load, held equal across architectures; the deployed first eigenfrequency is an equally defensible driver, and little hangs on the choice provided every architecture is sized against the same one.

    \item \textbf{Where the common-member assumption breaks.} The families do not all reduce to a monolithic member, and each departs differently. Trusses and masts are assemblies rather than single members, typically larger in envelope but lighter for a given stiffness; the driven linkage has members of unequal length and is only comparable once fully deployed; telescoping booms have progressively smaller segments, for which the telescopic-boom work of Wilting gives worked formulas that are adopted rather than re-derived. A realistic thickness is a requirement of the method rather than a refinement: with the members left dimensionless, every architecture stows to no volume at all and the packing comparison loses its meaning. A member of constant thickness is in any case not the optimal design, since tapering from the base is available to an articulated boom and comes automatically to a telescopic one, so the residual uncertainty this leaves is stated alongside the comparison rather than claimed away.

    \item \textbf{The packing-ratio boundary.} The figure is the volumetric ratio of the M2 mechanism stowed to the same mechanism deployed, bounded where it mounts to the rest of the spacecraft and running to its furthest extent: not the individual member, not the satellite, and not other deployables. Each architecture reports two envelopes, the minimum envelope around the wireframe and the box it fits into, giving a range rather than one number, with a qualitative note on how practical the resulting shape is. The box rather than the cylinder is the envelope of interest because a microsatellite of this class reaches orbit as a rideshare, where the accommodation is a rectangular volume on an interface ring, the outermost component decides which volume slot is occupied, and cost follows volume rather than mass. The volume the baffle is forced to occupy is recorded as a separate figure beside the ratio, never inside it, because a concept that flares at its base when stowed, or that swings outside the bus, pushes the baffle outward and so enlarges the stowed volume at system level. Because the baffle and the M2 mechanism are mechanically decoupled and deploy in sequence, a mechanism is permitted to swing outside the baffle and settle back inside it, so this system-level penalty is the only place where such a path shows up as a cost.

    \item \textbf{The optical parameters, carried as ranges.} The instrument design does not exist yet, and an uncertainty range will remain when it does, so these enter as ranges rather than as figures. The primary aperture is nominally \SI{30}{\centi\meter}, bounded \SI{20}{\centi\meter} to \SI{40}{\centi\meter}. The M1 to M2 separation is ideally equal to the aperture and is held below \SI{60}{\centi\meter} in any case, a ratio of two to one, which bounds the sweep of \Cref{sec:phase1} rather than leaving the separation free. An on-axis three-mirror configuration is the working assumption. The mass of M2 together with a possible fine adjustment stage is carried as an assumed range with its basis stated, not as a single value.
\end{itemize}

Two quantities are deliberately excluded. Aerodynamic drag moments on the deployed structure are not modelled, because the baffle shields both M2 and its support structure and carries that load. M2 placement error is not an output of this phase: the phase-1 models return geometry, the two stowed envelopes and the deployed envelope, member lengths, a first-order mass estimate and the sizing stiffness, and the error itself is built in \Cref{sec:phase2}.

The output of this activity is a written input-output specification. It exists so that the modelling that follows is aimed at something, rather than the parameters being settled retrospectively by whatever the first model happened to produce.

The output of this activity is a written input-output specification. It exists so that the modelling that follows is aimed at something, rather than the parameters being settled retrospectively by whatever the first model happened to produce.

\activity{models}{build the parametric models} Each architecture carried forward in \Cref{sec:families-assessed} is written as a parametric model in Python, with aperture and M1 to M2 separation as inputs rather than fixed dimensions, to the specification of \cref{act:inputs}.

\activity{sweep}{sweep the solution space} The models are run across the range of M1 to M2 separation and the aperture range, giving one curve per architecture on one common set of assumptions. The switch point is the separation at which the ranking between architectures inverts. The sweep is run against both a relative measure, separation divided by mirror diameter, and an absolute separation, as sub-question 2 requires. The sweep returns the two stowed envelopes and the deployed envelope, the baffle displacement recorded separately, member lengths, first-order mass estimate, and whether the geometry closes at all.

\activity{tradeoff}{the concept trade-off} Candidates that fail a first-order feasibility check, on geometric availability at the chosen separation or on stowed volume, are closed out with that reason recorded. The rest are scored against the criteria of \cref{act:criteria} on the sweep results, using the Analytic Hierarchy Process with a graphical trade-off as a cross-check, following the predecessor theses.

\activity{hinge}{the hinge and locking trade} \Cref{sec:properties} establishes that what defines the deployed position is set at the hinge rather than by the family, so the hinge is traded separately once the architecture is chosen. Its criteria are the three cases of \Cref{sec:properties}: a hard stop or latch, a driven position with feedback, or a shape found from stored elastic energy. This trade is decided from the literature.

% This trade is decided from the literature, with one argument carried over rather than inherited. The compliant rolling-contact hinge of the Delft heritage was adopted to keep the deployed structure from being overdetermined, yet whether a locked hinge really warps enough to matter has never been simulated, nor whether a compliant hinge introduces effects of its own under transient or microvibration input. That question is stated here and left open, because it is the same question as the latch count of \Cref{sec:deferred} and is answered there rather than in the hinge trade itself.

\activity{cad}{first model of the selected concept} The selected design is drawn as a CAD model: clearances, the swept arc, hinge and interface layout, and a packing ratio retrieved from the model rather than from the wireframe. Material and layup choices are made here, including resistance to atomic-oxygen exposure at VLEO, which is a material-selection question rather than a result of the thermoelastic analysis. This is the working model expected at the midterm.

\activity{phasetwoinputs}{gather and define the phase-2 inputs} A preparatory activity, placed deliberately before the midterm. The thermal environment for the \SI{300}{\kilo\meter} baseline and the assumed microvibration input spectrum are established here, each with the short piece of literature study it needs and an explicit decision on how it will be modelled. Gathering and defining inputs routinely costs more time than the analysis they feed, which is short work once the approach is settled, so doing it before the midterm means that even if it is unfinished the size of what remains is known, and any input that is simply missing has surfaced while there is still time to work around it.

\section{Phase 2, to the green light: analysing the selected design}
\label{sec:phase2}
\phasestart{2}

\activity{esatan}{ESATAN thermal cases} Orbital thermal cases for the \SI{300}{\kilo\meter} VLEO baseline, including eclipse cycling, with the cooled detector carried as a fixed boundary condition. The output is the temperature field and the gradients across the deployed structure over an orbit.

\activity{fem}{thermoelastic and modal finite-element analysis} The temperature fields drive a thermoelastic analysis of M2 pose drift, and a modal analysis gives the deployed first eigenfrequency for launch survival and for the reaction-wheel microvibration response. Material and layup choices are set in activity 1.7 and are taken as given here.

\activity{budget}{the error budget} A root-sum-square error budget propagates the contributors to M2 pose error: bare deployment repeatability, hinge or latch hysteresis, thermoelastic drift per orbit, manufacturing and alignment tolerances, and microvibration. This budget splits the sub-micron target into a number for each contributor, so every part of the design can be checked.

\activity{passiveactive}{the passive-versus-active assessment} The residual remaining after all passive measures is compared against the three degree-of-freedom correction budget: piston plus tip and tilt at the three arm bases. This comparison answers sub-question 3, and states the stroke and resolution the fine stage would require, sized against published corrector stages \cite{cadiergues_mirror_2003, lu_design_2023}.

\section{Two decisions deferred to the midterm}
\label{sec:deferred}

% TODO (Jasper, 27 Aug 2026) DECISION NEEDED, and it applies to this whole section.
% His position: a question carried as "decision deferred to the midterm" is in practice a question dropped at the midterm, because by then the remaining time decides it, and anything not explicitly planned for is unlikely to happen. His estimate is that at the midterm the latch/overdefinition study will be judged no longer feasible. So each of the two items below has to be planned explicitly now, or let go on purpose. He was clear that this is Hugo's thesis and he is not saying yes or no; he asked for about a week of mapping before deciding. Counter-argument he accepted for the latch study: the thermoelastic analysis (2.2) and the microvibration response are already planned activities, so pushing an over-determined model through the same process gives a second answer to compare at little extra cost. He also noted the question may already be relevant at activity 1.6 (where it can be left open) and at 1.7, so it could be decided straight after the trade-off in 1.5 rather than at the midterm. If the latch study is dropped, a lot of work falls away and the breadboard may become realistic again.

\textbf{A breadboard model.} Building hardware is proposed as a midterm decision, gated on the trade-off outcome, rather than committed to now. Its purpose would be to validate the CAD design physically and to run physical tests: that the mechanism folds and deploys as drawn, that the clearances hold, and that the deployment repeats. It is not an alternative route to the simulation results, and it would not measure at the accuracy that matters. The heritage repeatability campaign on comparable hardware reached 19 to \SI{25}{\micro\meter}, on one axis, in ambient conditions, which is more than an order of magnitude above the PLASMO target.

\textbf{The latch-count and overdefinition question.} Whether an overdefined structure with many latches warps under thermal stress, and whether latching only the middle joints while leaving the others free prevents that warping at the cost of vibration susceptibility, is an open question. It is assumed in earlier work and has not been investigated. It is a second research thread rather than part of the main line. The literature on it is checked during phase 1, and because the thermoelastic and modal analysis of activity 2.2 is already planned, running an over-determined variant of the same model through that activity would give a comparison at limited extra cost.

\section{Mapping of questions to activities}
\label{sec:mapping}

\begin{table}[H]
    \centering
    \caption{Mapping of research questions to research activities.}
    \label{tab:mapping}
    \small
    \begin{tabularx}{\textwidth}{p{4.0cm}p{2.4cm}X}
        \toprule
        Question & Primary activity & Output that answers it \\
        \midrule
        
        Sub-question 1: concepts and mechanisms & 1.1 to 1.6 & A trade-off on a common basis against the criteria fixed in activity 1.1, with one architecture selected and every discard reasoned, and a separate hinge and locking trade on the criteria of \Cref{sec:properties} \\
        \addlinespace
        
        Sub-question 2: validity across the optical space & 1.2, 1.3, 1.4 & The switch point in both relative and absolute form, and a statement on whether one concept spans both regimes \\
        \addlinespace
        
        Sub-question 3: the passive-active split & 1.8, 2.1 to 2.4 & An RSS error budget in five degrees of freedom, the residual compared against what each actuation location can reach, and a statement of the stroke and resolution a fine stage would need \\
        \addlinespace
        
        Main question & All of the above & A design with its packing ratio, its accuracy and stability budget, and its stiffness, stated against the open optical configuration \\
        
        \bottomrule
    \end{tabularx}
\end{table}

\section{Verification and working rules}
\label{sec:verification}

\begin{itemize}
    \item Every model is checked by hand, and every simulation campaign is treated as a test campaign, with objectives, conditions, an error analysis, and a requirement check at the end.
    
    \item Requirements are stated early and numbered, and every design chapter and simulation closes by checking against them. Because no deployment-accuracy requirement exists, the derived budget is itself traceable and is fed back to the design report.
    \item Requirements are stated early and numbered, and every design chapter and simulation closes by checking against them. Where a requirement has not yet been set by the instrument design, the budget derived in activity 2.3 stands in its place, is traceable to its inputs and assumptions, and is fed back to the design report; each such placeholder is replaced, throughout the document, as the requirement lands.
    
    \item Criteria are defined before concepts are scored, and discarded options are recorded with their reason.
    
    \item Writing runs alongside the work rather than at the end, so that a thesis outline and drafted chapters exist by the midterm.
    
    \item The design is kept generic under concurrent engineering, with open TBDs and TBCs carried explicitly rather than resolved by assumption.
\end{itemize}

\section{Resources, tools and risks}
\label{sec:risks}

\begin{table}[H]
    \centering
    \caption{Risk register for the research plan.}
    \label{tab:risks}
    \small
    \begin{tabularx}{\textwidth}{XXX}
        \toprule
        Risk or dependency & Impact & Mitigation \\
        \midrule
        
        Optical configuration still open (aperture, on-axis or off-axis, M1 to M2 distance) & Could invalidate a fixed design & Parametric models by construction; the switch-point sweep converts the uncertainty into a result. \\
        \addlinespace
        
        No deployment-accuracy requirement defined & No target to design against & Derive the budget within the thesis and feed it back to the design report; carry it as an explicit assumption until confirmed. \\
        \addlinespace
        
        Every architecture may exceed the error budget & Accuracy stops discriminating between architectures, and active compensation becomes necessary whatever is chosen & The budget of activity 2.3 is built early enough to be checked against the phase-1 results; if it is exceeded by all candidates, the criterion is dropped from the trade-off and the remaining actuator question is range rather than selection. \\
        \addlinespace
        
        The truss and mast architecture has no published thermal evidence in its family & An architecture could be selected on an untested assumption & Thermal behaviour enters the selection through its architectural proxies, the number of interfaces in the load path and the symmetry of the layout, and the absence of direct evidence is recorded as a stated assumption. The first direct evidence arrives with the thermoelastic analysis of activity 2.2, which is the gate on the assumption if this family is selected. \\
        \addlinespace
        
        Scope growth from the latch-warping thread & Two research threads in one thesis & Bounded by reusing the activity 2.2 model rather than building a separate study; the decision to include it is taken in phase 1, not left to the midterm. \\
        
        \bottomrule
    \end{tabularx}
\end{table}
